\chapter{Memoria descriptiva}
\thispagestyle{memoria}\pagestyle{memoria}

\section{Objeto}
El presente documento desarrolla el anteproyecto academico de instalaciones electricas en baja tension para un grifo de combustibles liquidos. El trabajo es nuevo y pertenece a Aquiles Taylor Ramos Yapo. El DXF recibido se usa exclusivamente para rescatar la implantacion, los ambientes, las islas, los surtidores y la zona de tanques; no se presenta como un proyecto electrico previo ni se modifica su fuente original.

Miguel Mamani Chuquicallata se identifica como propietario porque ese nombre aparece en la documentacion tecnica de referencia facilitada por la DREM. Esta transcripcion no equivale a una verificacion registral ni a una aprobacion de la DREM sobre el presente anteproyecto.

\section{Ubicacion y arquitectura}
El establecimiento se ubica en el predio rustico Reumita, parcelas B-8 y B-9, Comunidad Campesina San Francisco de Buenavista, carretera Juliaca--Puno, distrito de Caracoto, provincia de San Roman, Puno. La fuente muestra 5983.00 m\(^2\) de terreno, 2167.00 m\(^2\) destinados al proyecto, marquesina de 366.75 m\(^2\), zona de tanques de 108.24 m\(^2\), tres islas, seis surtidores y un edificio administrativo de tres niveles.

Se consideran cuatro tanques: dos de Diesel B5 S-50 de 9200 galones cada uno, uno de Gasohol Regular de 6580 galones y uno de Gasohol Premium de 6580 galones. GLP y GNV quedan expresamente fuera del alcance.

\section{Alcance electrico}
El anteproyecto incluye acometida de diseño, medicion, TGE, TDF, tableros administrativos, tablero de emergencia, ATS, grupo electrogeno, UPS, alumbrado normal y critico, tomacorrientes, bombas sumergibles, surtidores, control ATG, compresor, bombas de servicio, puesta a tierra, equipotencialidad, proteccion contra sobretensiones, proteccion contra el rayo y propuesta academica de areas peligrosas.

No incluye factibilidad de concesionaria, calculo definitivo de cortocircuito/selectividad, obra civil, instalaciones de GLP/GNV, proyecto sanitario/contra incendio completo ni habilitacion profesional. Estos componentes requieren documentos o especialidades que no fueron entregados.

\section{Suministro y continuidad}
Se adopta 3~x~380/220 V, 60 Hz, 3F+N+PE, suministro directo en baja tension y 50 kVA de capacidad de diseño. La propuesta queda condicionada a la factibilidad de Electro Puno y a la corriente de cortocircuito disponible. El neutro y el conductor de proteccion se mantienen separados aguas abajo del origen.

Para continuidad se propone un grupo de 37.5 kVA standby, ATS de cuatro polos y UPS senoidales para combustible/control y para POS--CCTV--comunicaciones. El respaldo es un criterio de ingenieria del anteproyecto, no una obligacion general atribuida sin matiz a todos los grifos.
